% ===================================================================
%  ТАБЛИЦЯ № 3 — Дитинство і Юність.
%  Traduction 2026 d'apres texto/io, controlee sur texto/fr, texto/en
%  et texto/es. Consigne : voir texto/uk/10-tablycia-01.tex.
%
%  Deux scenes, et l'ido subdivise beaucoup plus que le francais :
%  huit des intertitres n'ont pas de vis-a-vis francais. L'ukrainien
%  suit l'ido, comme les autres traductions.
%
%  LE FAC-SIMILE OUBLIE DE METTRE « julio » EN GRAS au bloc c2-01-1,
%  la ou il l'a mis a « junio » et a « augusto ». La source ne bouge
%  pas ; la correction est declaree dans gravuri/korekti.json, et
%  cette colonne y ajoute sa ligne — « (липня або ».
% ===================================================================

%%K t03-apar-1 apar
\VUsaut{3.0mm}
\VUtitre{15.0pt}{60}{ТАБЛИЦЯ № 3}
\VUsaut{1.6mm}
\VUtitre{11.4pt}{0}{Дитинство і Юність.}
\VUsaut{3.4mm}

%%K t03-apar-2 apar
(Таблиці 3 і 4 складено так, щоб дати в чотирьох \VUgras{родинних сценах}
основний запас слів про \VUgras{погоду}, \VUgras{вік}, \VUgras{родину},
\VUgras{одяг} і \VUgras{здоров'я}.)
\VUsaut{4.0mm}

%%K t03-tit-1 sub
\VUcentre{12.0pt}{0}{\textit{Перша сцена.}}
\VUsaut{3.0mm}

%%K t03-tit-2 sub
\VUpk{10.2pt}{40}{Хрестини на селі.}
\VUsaut{2.4mm}

%%K t03-tit-3 sub
\VUpk{10.2pt}{40}{Весна.}
\VUsaut{3.0mm}

%%K t03-c1-01-1 p
1. --- Ми у \VUgras{весні}, у місяці \VUgras{березні}, \VUgras{квітні} чи
\VUgras{травні}, в один із днів \VUgras{тижня} --- у \VUgras{понеділок},
\VUgras{вівторок} чи \VUgras{середу}. \VUgras{Десята година} ранку.

%%K t03-c1-02-1 p
2. --- \VUgras{Погода} чудова, \VUgras{повітря} чисте. Світить сонце. Кілька
маленьких \VUgras{хмаринок} \textsuperscript{(2)} підіймаються в синьому
\VUgras{небі} \textsuperscript{(1)}.

%%K t03-c1-03-1 p
3. --- На \VUgras{деревах} майже немає \VUgras{листя}. Стрункі \VUgras{тополі}
\textsuperscript{(3)} і високий \VUgras{платан} \textsuperscript{(17)} укриті
поки що тільки бруньками.

%%K t03-c1-04-1 p
4. --- \VUgras{Ластівки} \textsuperscript{(4)} вернулися і з радісними криками
кружляють навколо \VUgras{шпиля} \textsuperscript{(7)} \VUgras{дзвіниці}
\textsuperscript{(5)}, що вінчає сільську \VUgras{церкву} \textsuperscript{(6)}.

%%K t03-tit-4 sub
\VUsaut{3.4mm}
\VUpk{10.2pt}{40}{Церква.}
\VUsaut{3.0mm}

%%K t03-c1-05-1 p
5. --- Дуже проста ця церква з великим \VUgras{порталом}
\textsuperscript{(13)} і великим \VUgras{годинником} \textsuperscript{(8)}. Мала
\VUgras{стрілка} \textsuperscript{(9)} стоїть на цифрі X: вона показує
\VUgras{години}. \VUgras{Велика} \textsuperscript{(10)} стоїть на XII (полудень);
вона показує \VUgras{хвилини}.

%%K t03-c1-06-1 p
6. --- На дзвіниці весело дзвонять \VUgras{дзвони} \textsuperscript{(11)}. На
вершку даху стоїть невеликий кам'яний \VUgras{хрест} \textsuperscript{(12)}.

%%K t03-c1-07-1 p
7. --- У церкви не тільки великий \VUgras{портал} \textsuperscript{(13)}, є й
маленькі бічні двері. Через ці двері \VUgras{священник} \textsuperscript{(15)} у
\VUgras{сутані} виходить із \VUgras{ризниці} \textsuperscript{(14)} і прямує до
\VUgras{дому священника} \textsuperscript{(19)}.

%%K t03-c1-08-1 p
8. --- Поруч із ним --- \VUgras{молочарка} \textsuperscript{(24)} зі своїм
\VUgras{ослом} \textsuperscript{(23)}. Вона вертається з міста, куди возила дітям
своє добре молоко.

%%K t03-tit-5 sub
\VUsaut{4.0mm}
\VUpk{10.2pt}{40}{Садок.}
\VUsaut{3.4mm}

%%K t03-c1-09-1 p
9. --- Пан Сірий (осел) дуже вдоволений цією ранковою прогулянкою. Він з
насолодою вдихає повітря, напахчене \VUgras{цвітом} \textsuperscript{(20)}, який
укриває \VUgras{фруктові дерева} \textsuperscript{(18)} сусіднього
\VUgras{саду}. Усі в рожевому й білому ці \VUgras{персикові дерева}
\textsuperscript{(18)} і \VUgras{абрикосові дерева} \textsuperscript{(19)}, і
вони обіцяють добрий урожай.

%%K t03-c1-10-1 p
10. --- А тим часом маленькі \VUgras{бджоли} \textsuperscript{(22)} з
\VUgras{вулика} \textsuperscript{(21)} летять туди збирати свій солодкий
\VUgras{мед}, гудучи за роботою.

%%K t03-c1-11-1 p
11. --- Але що сталося? Ось біжать сільські діти й розганяють настрашених
\VUgras{гусей} \textsuperscript{(25)}. Це хода виходить із церкви після
\VUgras{хрестин} нотаревого сина.

%%K t03-c1-12-1 p
12. --- Малий Яків у своїй \VUgras{колисці} \textsuperscript{(26)} прокидається
від веселого дзвону дзвонів, а його сестра Магдалина, якій теж хочеться
подивитися на ходу, просить матір зачесати й одягти її на порозі.

%%K t03-c1-13-1 p
13. --- Мати погоджується. Вона надіває свою \VUgras{хустку}
\textsuperscript{(28)}, свою \VUgras{кофту} \textsuperscript{(29)} і свій
\VUgras{фартух} \textsuperscript{(30)} і сідає на низький стілець перед дверима.
На Магдалині тільки \VUgras{спідниця} \textsuperscript{(31)} і \VUgras{ліфчик}
\textsuperscript{(32)}.

%%K t03-c1-14-1 p
14. --- Яка вона щаслива! Вона забуває свої улюблені квіти, свої
\VUgras{гвоздики} \textsuperscript{(34)}, на які сідають \VUgras{метелики}
\textsuperscript{(35)}, свої \VUgras{тюльпани} \textsuperscript{(36)}, свої
\VUgras{конвалії} \textsuperscript{(37)} у \VUgras{горщиках}
\textsuperscript{(38)} на \VUgras{мурі} \textsuperscript{(33)} і свої
\VUgras{троянди} \textsuperscript{(39)}, які утворюють такий гарний живопліт.
Вона задивляється на багаті вбрання дам у ході.

%%K t03-c1-15-1 p
15. --- Сільські хлопці мало дивляться на вбрання. Вони кидаються вперед і
штовхають одне одного, щоб дістати \VUgras{цукерки} \textsuperscript{(55)} або
\VUgras{монети} \textsuperscript{(52)}. Один із них плаче
\textsuperscript{(40)}.

%%K t03-c1-16-1 p
16. --- Другий наступив на лапу \VUgras{собаці} \textsuperscript{(42)}, і той із
скавучанням тікає й обертає на втечу \VUgras{курку} \textsuperscript{(43)} з її
\VUgras{курчатами} \textsuperscript{(44)}.

%%K t03-tit-6 sub
\VUsaut{5.0mm}
\VUpk{10.2pt}{40}{Хрестинна хода.}
\VUsaut{4.0mm}

%%K t03-c1-17-1 p
17. --- На чолі ходи йде \VUgras{годувальниця} \textsuperscript{(45)}. На ній
гарний \VUgras{чепець} \textsuperscript{(46)} з довгими стрічками. Вона несе
\VUgras{новонародженого} \textsuperscript{(47)}, усього в \VUgras{мереживі}
\textsuperscript{(48)}.

%%K t03-c1-18-1 p
18. --- За годувальницею йдуть \VUgras{хрещений батько} \textsuperscript{(49)} і
\VUgras{хрещена мати} \textsuperscript{(53)}. Вони роздають цукерки й монети. У
хрещеного батька \VUgras{рукавички} \textsuperscript{(51)} і \VUgras{циліндр}
\textsuperscript{(50)}. Хрещена мати тримає в руці гарний \VUgras{букет}
\textsuperscript{(54)}.

%%K t03-c1-19-1 p
19. --- За ними видно \VUgras{дядька} \textsuperscript{(56)} і \VUgras{тітку}
\textsuperscript{(58)} дитини. На тітці зграбний \VUgras{капелюх}
\textsuperscript{(59)}, на дядькові чорний \VUgras{сюртук}
\textsuperscript{(57)} і білий жилет. Далі йдуть \VUgras{двоюрідний брат}
\textsuperscript{(60)} і \VUgras{двоюрідна сестра} \textsuperscript{(61)}.
Остання веде за руку \VUgras{сестру} \textsuperscript{(62)} новонародженого,
малу Берту.

%%K t03-c1-20-1 p
20. --- Другий \VUgras{брат} \textsuperscript{(63)} Берти йде позаду. Він подає
руку \VUgras{родичці} \textsuperscript{(64)}. За ними йдуть \VUgras{друзі}
\textsuperscript{(65)} родини.

%%K t03-tit-7 sub
\VUsaut{4.6mm}
\VUpk{10.2pt}{40}{Роззяви.}
\VUsaut{3.4mm}

%%K t03-c1-21-1 p
21. --- Обіч ходи --- \VUgras{жебрак} \textsuperscript{(66)}, що спирається на
свою \VUgras{милицю} \textsuperscript{(67)}. Він просить милостині.

%%K t03-c1-22-1 p
22. --- По другий бік --- дуже \VUgras{чемний} \textsuperscript{(68)} хлопець.
Він підіймає шапку і каже \VUgras{«доброго дня»} тим, кого знає.

%%K t03-c1-23-1 p
23. --- Селяни вийшли з домів подивитися на хрестинну ходу. Ми бачимо
\VUgras{селянина} \textsuperscript{(69)} в \VUgras{паперовому ковпаку}, а поруч
із ним \VUgras{селянку} \textsuperscript{(70)}. Тоді \VUgras{стару жінку}
\textsuperscript{(71)}, що спирається на свою \VUgras{палицю}
\textsuperscript{(72)}. Нарешті \VUgras{мірошника} \textsuperscript{(73)} в
широкому \VUgras{повстяному капелюсі} \textsuperscript{(74)} і молоду селянку в
\VUgras{дерев'яних черевиках} \textsuperscript{(75)}, яка вчить ходити свого
малюка.

%%K t03-c1-24-1 p
24. --- Картина дуже жвава.
\VUsaut{4.0mm}

%%K t03-tit-8 sub
\VUcentre{12.0pt}{0}{\textit{Друга сцена.}}
\VUsaut{3.0mm}

%%K t03-tit-9 sub
\VUpk{10.2pt}{40}{Народне свято.}
\VUsaut{2.4mm}

%%K t03-tit-10 sub
\VUpk{10.2pt}{40}{Водні змагання і перегони.}
\VUsaut{3.0mm}

%%K t03-c2-01-1 p
1. --- Прийшло \VUgras{літо}. Палюче сонце місяця \VUgras{червня} (липня або
\VUgras{серпня}) вабить молодь купатися і кататися на човнах.

%%K t03-c2-02-1 p
2. --- На правому березі річки веслярі роздягаються, щоб узяти участь у водних
змаганнях і в перегонах.

%%K t03-c2-03-1 p
3. --- Один із них скидає \VUgras{черевики} \textsuperscript{(1)}; він
розшнуровує (розстібає) їх. Двоє його товаришів уже готові. На них тепер тільки
\VUgras{фуфайка} \textsuperscript{(2)} і \VUgras{плавки} \textsuperscript{(3)}.
Вони балакають, чекаючи на товаришів. Вони говорять про ярмарок і про свято на
тому березі.

%%K t03-c2-04-1 p
4. --- На землі поруч із ними лежить \VUgras{одяг} їхніх товаришів:
\VUgras{пара} \VUgras{чобіт} \textsuperscript{(4)}, \VUgras{туфлі}
\textsuperscript{(5)}, \VUgras{шкарпетки} \textsuperscript{(6)},
\VUgras{повстяний} \textsuperscript{(7)} і \VUgras{солом'яний}
\textsuperscript{(8)} капелюхи та \VUgras{краватка} \textsuperscript{(9)}.

%%K t03-c2-05-1 p
5. --- Один із юнаків охайно склав свій \VUgras{піджак} \textsuperscript{(10)}.
Він кладе його на рушник, у який його сусід зараз загорне обидва їхні
\VUgras{костюми}.

%%K t03-c2-06-1 p
6. --- Шостий хлопець запізнюється. На голові в нього ще \VUgras{кашкет}
\textsuperscript{(11)}. Він не скинув ні \VUgras{сорочки} \textsuperscript{(12)},
ні \VUgras{комірця} \textsuperscript{(13)}, ні \VUgras{манжет}
\textsuperscript{(14)}. Він тримає в руці свій \VUgras{жилет}
\textsuperscript{(17)} і все ще в \VUgras{штанях} \textsuperscript{(16)} і
\VUgras{шлейках} \textsuperscript{(15)}. На наступний заїзд він напевно не
встигне.

%%K t03-c2-07-1 p
7. --- Обіч юнаків стежка, обросла \VUgras{будяками} \textsuperscript{(18)},
спускається до води, де старий Яків готує серед \VUgras{очерету}
\textsuperscript{(19)} \VUgras{феєрверк} \textsuperscript{(20)}, який пустять
увечері.

%%K t03-tit-11 sub
\VUsaut{4.4mm}
\VUpk{10.2pt}{40}{Перегони.}
\VUsaut{3.4mm}

%%K t03-c2-08-1 p
8. --- На річці кілька дам, сховавшись під парасольками, катаються в
\VUgras{човні} \textsuperscript{(21)}. Вони стежать за перегонами, а вони дуже
цікаві. У кожній \VUgras{гічці} \textsuperscript{(22)} четверо \VUgras{веслярів}
\textsuperscript{(24)} веслують щосили, а \VUgras{стерновий}
\textsuperscript{(26)} тримає \VUgras{стерно} \textsuperscript{(25)} і править
крихким суденцем. \VUgras{Весла} \textsuperscript{(23)} мірно б'ють по воді, і
легкі човни мчать із запаморочливою швидкістю.

%%K t03-c2-09-1 p
9. --- Кілька юнаків змагаються біля бика мосту за приз \VUgras{щогли з призом}
\textsuperscript{(28)}. Один із них обережно посувається по пружкому й слизькому
жердині, щоб дістати \VUgras{прапорець} \textsuperscript{(29)}, закріплений на
кінці. Його невдатний \VUgras{суперник} \textsuperscript{(30)} упав у воду. Він
пливе назад до берега.

%%K t03-c2-10-1 p
10. --- Старші хлопці \VUgras{б'ються на воді} \textsuperscript{(32)} біля
\VUgras{набережної} \textsuperscript{(33)}. Великий \VUgras{натовп}
\textsuperscript{(31)} дивиться на всі ці ігри, спершись на \VUgras{поручні}
\VUgras{мосту} \textsuperscript{(27)} і тиснучись уздовж \VUgras{берега}. Інші
роззяви, щоправда, обертаються, щоб поглянути на гру зі \VUgras{стовпом}
\textsuperscript{(45)} або помилуватися \VUgras{ходою мерії}, що йде на
\VUgras{роздачу нагород}.

%%K t03-c2-11-1 p
11. --- Спершу біжать кілька \VUgras{хлопчаків} \textsuperscript{(35)} поперед
\VUgras{музикантів} \textsuperscript{(36)}; тоді йдуть \VUgras{мер}
\textsuperscript{(37)}, його помічники і \VUgras{міська рада} містечка.
Останніми крокують \VUgras{пожежники} \textsuperscript{(38)}.

%%K t03-tit-12 sub
\VUsaut{5.0mm}
\VUpk{10.2pt}{40}{Ярмарок.}
\VUsaut{3.6mm}

%%K t03-c2-12-1 p
12. --- На \VUgras{ярмарковій площі} \textsuperscript{(39)} великий натовп. Люди
приходять подивитися на різні \VUgras{балагани}: на \VUgras{дерев'яних коників}
\textsuperscript{(40)}, на \VUgras{цирк} \textsuperscript{(41)}, де вистава вже
почалася; на \VUgras{народний бал} \textsuperscript{(42)}, де молоді люди й
дівчата містечка віддаються втісі \VUgras{танцювати}.

%%K t03-c2-13-1 p
13. --- У глибині видно \VUgras{арену} \textsuperscript{(44)}, де відважний
іспанський \VUgras{матадор} б'ється з розлюченим \VUgras{биком}
\textsuperscript{(45)}; тоді \VUgras{російські гори} \textsuperscript{(49)} і
\VUgras{тир} \textsuperscript{(46)}, де вправляються в поводженні з
\VUgras{вогнепальною зброєю} і в стрільбі по \VUgras{мішені}.

%%K t03-c2-14-1 p
14. --- Поруч --- \VUgras{балаганний театр} \textsuperscript{(47)}, де грають
\VUgras{комедію} і \VUgras{драму}. Зовсім близько стоїть балаган
\VUgras{велетня} \textsuperscript{(48)} і \VUgras{карлика}. \VUgras{Зріст}
першого --- два метри п'ятнадцять; другий ледве з дитину.

%%K t03-c2-15-1 p
15. --- Нарешті, ось великий \VUgras{звіринець} \textsuperscript{(50)} з його
\VUgras{дикими звірами}, з \VUgras{тигром} \textsuperscript{(51)} і його
могутніми \VUgras{пазурами}, з \VUgras{левом} \textsuperscript{(52)} і його
густою \VUgras{гривою}, з двома \VUgras{жирафами} \textsuperscript{(53)} і
\VUgras{слоном} \textsuperscript{(54)}, який так спритно орудує своїм
\VUgras{хоботом}.

%%K t03-c2-16-1 p
16. --- \VUgras{Приборкувач} \textsuperscript{(55)}, який навчає цих звірів,
стоїть біля дверей балагана.
\VUsaut{6.0mm}
\VUfilet{22mm}
